% ---------------------------------------------------------------------------
% OWNER: Advaya Bhardwaj - what students feel (supervised classification).
% Follow TEAM_INSTRUCTIONS.md, Part B. Same structure as Section 4 (Tushaar):
% copy its style. Figures are already wired below: once your script saves
% images/M2_*.png with these names they appear automatically.
% Replace every \todo box; delete this comment block when done.
% ---------------------------------------------------------------------------
\section{Advaya Bhardwaj: Dropout Risk (Classification)}
\label{sec:m2}

\paragraph{Driving question.}
\todo{One paragraph. Suggested: A1 (Appendix, Fig.~21) found that the
dropout-risk flag is a hard rule on stress and motivation. Can a classifier
predict dropout risk, and once the rule's own inputs are removed, does anything
about a student still predict it?}

\paragraph{Method summary.}
\todo{One paragraph: logistic regression, decision tree, random forest
(scikit-learn~\cite{pedregosa2011}); class weights for the 2\% positive rate;
metrics are PR-AUC, ROC-AUC and F1 on the held-out test split (accuracy is
useless at 98\% negatives); permutation importance. Fit on train only.}

\teamfig{\textwidth}{images/M2_workflow_unrolled.png}{Member~2's workflow,
unrolled into two iterations with the \emph{Feedback loop} edge labelled.
\todoi{finish caption}}{fig:m2-workflow}

\todo{Add a task/rationale table like Table~\ref{tab:m1-tasks}: one row per
figure with stage, task verb and why that chart type.}

\subsection{Iteration 1: all features}

\teamfig{0.75\textwidth}{images/M2_fig01_rule_heatmap.png}{\todoi{Caption: dropout
rate for every stress $\times$ motivation cell on the training split.}}{fig:m2-rule}

\teamfig{\textwidth}{images/M2_fig02_iter1_pr_roc.png}{\todoi{Caption: PR and ROC
curves for the three models, test split. Quote PR-AUC per model.}}{fig:m2-pr1}

\teamfig{\textwidth}{images/M2_fig03_iter1_tree.png}{\todoi{Caption: depth-2 tree
and importances. The tree should split on stress $>8.05$ then
motivation $\le3.5$.}}{fig:m2-tree}

\paragraph{Knowledge K1.}
\todo{Numbered inferences. Expect: tree and forest reach PR-AUC 1.00;
logistic regression only about 0.88 because a straight boundary cannot draw an
AND-rectangle; the tree recovers the exact thresholds. Conclusion: the label is
a definition, so the models read it back (label leakage).}

\subsection{Feedback loop: transforming the data}
\todo{Describe T1: drop \texttt{stress\_level} and \texttt{motivation\_level}
(\texttt{exam\_anxiety\_score} is already removed in Section~\ref{sec:a3-prep});
keep class weights. Say why each change follows from K1.}

\subsection{Iteration 2: leak-free features}

\teamfig{\textwidth}{images/M2_fig04_iter2_pr_before_after.png}{\todoi{Caption:
PR curves before and after removing the rule variables, with the 2\% base
rate as a reference line.}}{fig:m2-pr2}

\teamfig{\textwidth}{images/M2_fig05_iter2_importance.png}{\todoi{Caption:
permutation importance of the leak-free model, test split.}}{fig:m2-imp}

\teamfig{\textwidth}{images/M2_fig06_proxy_check.png}{\todoi{Caption: dropout
rate by prior-GPA band and by mental-health rating.}}{fig:m2-proxy}

\paragraph{Knowledge K2.}
\todo{Numbered inferences. Expect: PR-AUC falls to about 0.04--0.08 (base rate
0.02) and ROC-AUC to about 0.70. Check which features carry the remaining
signal; a quick check found prior GPA alone gives ROC 0.67, because prior GPA
correlates $r=0.27$ with motivation, and mental health steps down at stress 7
(A1 Fig.~23). So what is left is mostly proxies for the two rule variables. Say
what this means: the flag measures how students feel, not what they do.}
